\documentclass[12pt,a4paper]{report}
\usepackage[T1]{fontenc}
\usepackage[utf8]{inputenc}
\usepackage{mathptmx}         % Times New Roman (text + math)
\usepackage{amsmath,amsthm,amssymb,graphicx,hyperref,float}
\usepackage[left=3cm,right=2cm,top=2cm,bottom=2cm]{geometry}
\usepackage[romanian]{babel}
\usepackage{setspace}         % single line spacing
\usepackage{url}              % proper URL formatting in bibliography
\begin{document}
\singlespacing
\setlength{\parindent}{1.27cm}
\setlength{\parskip}{0pt}
\thispagestyle{empty}
\begin{center}
\begin{figure}[h!]
\vspace{-20pt}
\begin{center}
\includegraphics[width=100pt]{collegelogo.jpg}
\end{center}
\end{figure}


{\large{\bf Universitatea Tibiscus din Timișoara

Facultatea de Calculatoare și Informatică Aplicată}}

\vspace{120pt}
{\huge {\bf LUCRARE DE LICENŢĂ}}

\vspace{150pt}
\end{center}

{\large\noindent{\bf COORDONATOR:\hfill ABSOLVENT:}

\noindent Lect. Dr. Apostol Simona \hfill Groza Darius-Nicolae}

\vfill
\begin{center}
{\bf TIMIȘOARA

2026}
\end{center}
\newpage
\thispagestyle{empty}
\begin{center}
{\large{\bf Universitatea Tibiscus din Timișoara

Facultatea de Calculatoare și Informatică Aplicată}}

\vspace{120pt}
{\huge {\bf InfraMind -- Sistem pentru sincronizarea infrastructurilor IT distribuite}}

\vspace{150pt}
\end{center}

{\large\noindent{\bf COORDONATOR:\hfill ABSOLVENT:}

\noindent Lect. Dr. Apostol Simona \hfill Groza Darius}

\vfill
\begin{center}
{\bf TIMIȘOARA

2026}
\end{center}
\tableofcontents
\newpage
\chapter*{Introducere}
\addcontentsline{toc}{chapter}{Introducere}



\noindent În contextul actual al dezvoltării infrastructurilor IT, în care containerizarea și orchestrarea aplicațiilor distribuite joacă un rol central, gestionarea eficientă a clusterelor Kubernetes devine o necesitate tot mai prezentă, datorită numeroaselor beneficii ce vin cu acestea. Deși tehnologiile de tip Kubernetes sunt preponderente în lumea enterprise a companiilor, așa cum a fost confirmat și de un studiu realizat de Edge Delta, care arată că, în ultimii ani, peste 60\% din companiile IT au adoptat această tehnologie [2], utilizarea ei necesită un nivel crescut de cunoștințe, alăturat de multe abilități din domenii diferite, cum ar fi: Networking, Automatizare, Linux, Securitate, Containerizare, Orchestrare etc.

Astfel apare următoarea situație: nevoia clară de tot mai multe clustere Kubernetes, dar și problema: dificultatea operării acestora. InfraMind a fost creat ca un instrument dedicat automatizării și simplificării procesului de instalare și configurare a unui cluster Kubernetes, servind ca o soluție ce are ca scop ușurarea acestui proces pentru unul sau mai multe clustere. Soluția are la bază tehnologii moderne — K3s, Python/Paramiko, Flask, Bash, Docker — și nu depinde de niciun tool extern de tip configuration management.

Scopul principal al acestei aplicații este de a oferi o soluție completă și scalabilă pentru implementarea unui cluster K3s (o variantă lightweight a Kubernetes) pe infrastructuri variate. InfraMind automatizează atât procesul inițial de instalare, cât și configurările ulterioare necesare pentru gestionarea unui mediu containerizat, folosind exclusiv conexiuni SSH directe orchestrate prin Python.

Aplicația are la bază următoarele tehnologii principale:

\begin{itemize}
  \item \textbf{K3s} — distribuție Kubernetes lightweight, utilizată ca motor de orchestrare al clusterului.
  \item \textbf{Paramiko} — librărie Python pentru comunicare SSH, utilizată pentru orchestrarea instalării directe pe noduri.
  \item \textbf{Jinja2} — motor de template-uri Python folosit pentru generarea configurațiilor K3s per nod.
  \item \textbf{Flask} — micro-framework web Python folosit atât ca backend REST, cât și ca server al interfeței utilizatorului.
  \item \textbf{Docker} — platformă de containerizare utilizată pentru rularea portabilă a aplicației InfraMind.
\end{itemize}

În concluzie, InfraMind reprezintă o soluție modernă și eficientă pentru oricine dorește să implementeze și să gestioneze un cluster Kubernetes într-un mod automatizat, reducând complexitatea procesului și riscurile asociate erorilor de configurare manuală. Acest proiect demonstrează puterea automatizării în DevOps și importanța unor fluxuri de lucru bine structurate în gestionarea infrastructurilor moderne.

Lucrarea este organizată după cum urmează: Capitolul~\ref{cap:similare} prezintă soluțiile similare existente pe piață. Capitolul~\ref{cap:tehnologii} descrie tehnologiile utilizate în implementarea InfraMind. Capitolul~\ref{cap:cerinte} analizează cerințele funcționale și non-funcționale ale aplicației. Capitolul~\ref{cap:structura} detaliază arhitectura și structura internă a soluției.
\chapter{Soluții Similare}
\label{cap:similare}
\section{Soluții Enterprise}
Soluțiile de tip Enterprise oferă un număr mai mare de servicii și beneficii, însă la costuri mult mai mari, cu anumite limitări — cum ar fi tipul de infrastructură suportat — și nu toate pot fi la fel de versatile precum InfraMind.

\subsection{Red Hat OpenShift}

OpenShift este o platformă enterprise pentru dezvoltarea și livrarea aplicațiilor care include Kubernetes la bază, dar adaugă și caracteristici suplimentare pentru securitate, integrare CI/CD și gestionarea ușoară a aplicațiilor [16].

\begin{itemize}
  \item Management complet al clusterelor Kubernetes.
  \item Securitate îmbunătățită (Role-Based Access Control, OAuth etc.).
  \item Suport pentru aplicații de tip microservicii și integrare CI/CD.
  \item Tooling pentru dezvoltatori, inclusiv imagini Docker integrate și un catalog de aplicații.
  \item Deploy: instalabil atât pe infrastructură proprie (on-premise), cât și în cloud (AWS, Azure, GCP).
\end{itemize}

\subsection{Microsoft Azure Kubernetes Service (AKS)}

AKS este un serviciu complet gestionat de Kubernetes în Azure, care permite dezvoltatorilor să implementeze și să administreze ușor clustere Kubernetes [9]. Microsoft se ocupă de gestionarea infrastructurii, iar utilizatorii se pot concentra pe aplicațiile lor.

\begin{itemize}
  \item Gestionare completă a clusterului Kubernetes, inclusiv patch-uri și actualizări.
  \item Scalare automată și integrări cu alte servicii Azure (de exemplu, Azure Active Directory pentru autentificare).
  \item Integrări cu DevOps pentru CI/CD și monitorizare avansată.
  \item Deploy: exclusiv în cloud-ul Azure, gestionat complet de Microsoft.
\end{itemize}

\subsection{Rancher}

Rancher este o soluție open-source care oferă o platformă de gestionare a clusterelor Kubernetes pentru multiple infrastructuri (on-premise, AWS, Azure etc.) [14]. Este ideală pentru organizațiile care doresc un control mai mare asupra modului în care își gestionează infrastructura Kubernetes.

\begin{itemize}
  \item Gestionare unificată a mai multor clustere Kubernetes.
  \item Integrări pentru CI/CD, monitorizare și gestionarea securității.
  \item Multi-cloud și multi-cluster: suportă implementarea în mai multe medii (on-premise, AWS, Azure etc.).
  \item Deploy: poate fi implementat pe orice infrastructură, inclusiv medii cloud sau on-premise.
\end{itemize}

\subsection{Google Kubernetes Engine (GKE)}

GKE este un serviciu complet gestionat de Google pentru implementarea și gestionarea clusterelor Kubernetes [4]. Google, fiind creatorul Kubernetes, oferă un serviciu extrem de scalabil și bine integrat cu celelalte servicii Google Cloud.

\begin{itemize}
  \item Management complet al clusterelor Kubernetes cu upgrade-uri automate și patch-uri.
  \item Scalare automată pe baza cerințelor de trafic.
  \item Integrare completă cu alte servicii Google Cloud (BigQuery, Pub/Sub etc.).
  \item Deploy: exclusiv în Google Cloud, complet gestionat de Google.
\end{itemize}
\subsection{Beneficii și diferențe față de soluțiile Enterprise}

Un beneficiu direct al faptului că soluția InfraMind instalează clustere K3s folosind orice tip de infrastructură (fizică și virtuală) este flexibilitatea și portabilitatea maximă a implementării cluster-ului. Acest lucru înseamnă că utilizatorii pot implementa și gestiona clustere K3s indiferent dacă resursele sunt pe hardware fizic sau într-un mediu virtualizat, într-un cloud public sau privat, fără a fi limitați de tipul de infrastructură, spre deosebire de unele dintre soluțiile enterpise alternative care limitează utilizatorul la un singur tip de infrastructură, și posibil sub anumite reglementări/condiții.

Costuri reduse: Posibilitatea de a rula clustere K3s pe hardware existent, indiferent dacă este fizic sau virtualizat, poate ajuta organizațiile să evite cheltuielile suplimentare cu hardware dedicat.

Scalabilitate mai mare: Fiecare tip de infrastructură (fizică sau virtuală) poate fi scalat mai eficient, ceea ce oferă mai multe opțiuni pentru gestionarea resurselor.

Adaptabilitate: Organizațiile pot migra între infrastructuri fără întreruperi majore, păstrând aceleași configurații și procese de management ale clusterei Kubernetes.

Aceste aspecte fac soluția InfraMind foarte atractivă pentru utilizatorii care doresc să își optimizeze infrastructura IT și să își reducă complexitatea administrativă, beneficiind în același timp de avantajele K3s.

Cu toate aceste beneficii, InfraMind este limitat însă în partea high-end de raw power a lucrurilor, InfraMind fiind orientat spre volume de muncă mici până la medii, unde aplicațiile enterpirse încep de la volume medii spre cât mai mari.

\section{Soluții Open-source}

Aceste soluții se aseamănă mult cu aplicația InfraMind, folosind abordări minimale și econome.

\subsection{Kubeadm}

Kubeadm este o unealtă open-source creată de Kubernetes pentru a ajuta la instalarea și gestionarea unui cluster Kubernetes [7]. Este o soluție flexibilă care oferă un mod simplu de utilizat pentru a configura un cluster Kubernetes de la zero.

\begin{itemize}
  \item Configurare rapidă a unui cluster Kubernetes.
  \item Permite crearea unui cluster pe infrastructuri on-premise sau în cloud.
  \item Suportă funcționalități de upgrade, gestionare și inițializare a clusterelor.
\end{itemize}

\subsection{Minikube}

Minikube este o soluție open-source care permite crearea unui cluster Kubernetes local pe o singură mașină, fiind ideală pentru dezvoltatori și testare în medii de dezvoltare [10].

\begin{itemize}
  \item Instalare rapidă a unui cluster Kubernetes local pe mașini Windows și Linux.
  \item Suportă pluginuri și extensii.
  \item Permite testarea aplicațiilor în Kubernetes fără a necesita un cluster complet.
\end{itemize}

\subsection{Kind}

Kind este o unealtă care permite crearea de clustere Kubernetes folosind containere Docker [6]. Este ideală pentru dezvoltatori care doresc să testeze Kubernetes într-un mediu izolat pe mașinile lor locale.

\begin{itemize}
  \item Permite rularea de clustere Kubernetes complete în containere Docker.
  \item Ideal pentru testare rapidă și dezvoltare locală.
  \item Suportă crearea unui cluster cu mai multe noduri.
\end{itemize}

\subsection{Kubicorn}

Kubicorn este o soluție open-source care facilitează implementarea unui cluster Kubernetes în infrastructuri publice de cloud (AWS, Azure, Google Cloud) folosind o abordare declarativă.

\begin{itemize}
  \item Suportă implementarea automată de clustere Kubernetes în cloud.
  \item Permite implementarea în mai multe tipuri de infrastructuri.
  \item Focalizat pe automatizare și ușurința de gestionare.
\end{itemize}

\subsection{Beneficii și diferențe față de soluțiile Open-source}

Deși aceste aplicații au reprezentat, într-o oarecare măsură, o sursă de inspirație pentru InfraMind, acestea au un dezavantaj major față de InfraMind: lipsa unei interfețe grafice, care să ofere o experiență mai plăcută utilizatorului. De asemenea, folosirea acestor aplicații necesită, din nou, un nivel înalt de cunoștințe tehnice — lucru pe care dorim să îl evităm, scopul InfraMind fiind de a ușura cât mai mult procesul de instalare și configurare.

\chapter{Tehnologii Utilizate}
\label{cap:tehnologii}

Această secțiune prezintă principalele tehnologii și instrumente pe care se bazează InfraMind, motivând alegerea fiecăreia și descriind rolul ei în cadrul arhitecturii globale a soluției.

\section{K3s}

K3s este o distribuție Kubernetes certificată, creată de Rancher Labs (acum parte din SUSE), concepută special pentru medii cu resurse limitate, dispozitive IoT și scenarii edge [5]. Spre deosebire de Kubernetes clasic [8], K3s este ambalat ca un singur binar de mai puțin de 100 MB, eliminând dependențele externe și componentele considerate opționale în contextul unui deployment ușor.

Principalele caracteristici care au motivat alegerea K3s în cadrul InfraMind sunt:

\begin{itemize}
  \item \textbf{Consum redus de resurse} — K3s necesită semnificativ mai puțin RAM și CPU față de Kubernetes standard, ceea ce îl face potrivit atât pentru hardware fizic de capacitate medie, cât și pentru mașini virtuale cu resurse limitate.
  \item \textbf{Instalare simplificată} — Întregul proces de bootstrap al unui nod master sau worker se realizează printr-un singur script shell descărcat de la \texttt{get.k3s.io}, ușurând automatizarea prin SSH.
  \item \textbf{Compatibilitate completă cu ecosistemul Kubernetes} — K3s respectă API-ul standard Kubernetes, ceea ce înseamnă că orice manifest YAML sau unealtă (kubectl, Helm, Prometheus etc.) funcționează neschimbat.
  \item \textbf{Suport pentru clustere HA} — Prin utilizarea mecanismului de embedded etcd, K3s suportă configurații cu mai mulți masteri, așa cum este cazul topologiei implementate de InfraMind.
\end{itemize}

\section{Paramiko și orchestrarea SSH}

Paramiko este o librărie Python matură pentru comunicare SSH2, care oferă atât funcționalități client cât și server, bazată pe biblioteca criptografică \textit{cryptography} [11]. Spre deosebire de soluțiile tradiționale de tip configuration management (Ansible, Chef, Puppet), Paramiko permite integrarea directă a orchestrării SSH în codul Python al aplicației, fără dependențe externe de instalat pe mașina de control sau pe nodurile țintă.

În InfraMind, Paramiko constituie nucleul motorului de orchestrare și îndeplinește toate sarcinile de comunicare cu nodurile clusterului:

\begin{itemize}
  \item \textbf{Testarea conectivității SSH} — Înainte de orice instalare, fiecare nod este testat individual prin deschiderea unei sesiuni SSH cu cheia privată furnizată de utilizator.
  \item \textbf{Executarea comenzilor la distanță} — Comenzile de instalare și dezinstalare K3s, de instalare Docker și de gestionare a containerelor sunt executate pe noduri prin canale SSH non-blocante, cu citire în timp real a output-ului.
  \item \textbf{Transferul de fișiere} — Fișierele de configurare K3s generate prin Jinja2 sunt transferate pe noduri via \texttt{sudo tee}, fără a necesita un protocol suplimentar (SFTP/SCP).
  \item \textbf{Suport pentru mai multe formate de chei} — La deschiderea fiecărei conexiuni se încearcă succesiv formatele Ed25519, RSA, ECDSA și DSS, asigurând compatibilitatea cu orice infrastructură existentă.
\end{itemize}

Un aspect important al implementării este mecanismul de output non-blocant: conexiunile SSH folosesc canale fără blocare (\textit{non-blocking}), astfel încât liniile de output sunt yield-uite progresiv către frontend prin Server-Sent Events, fără a bloca procesul Flask.

Alegerea Paramiko față de un tool extern de tip configuration management a fost motivată de eliminarea oricărei dependențe externe pe mașina host, de integrarea nativă în codul Python și de controlul granular asupra comportamentului SSH (timeout-uri, abort, streaming de output).

\section{Jinja2}

Jinja2 este un motor de template-uri Python rapid și expresiv, creat de același echipă ca și Flask (Pallets Projects) [12]. Este utilizat extensiv în aplicații web pentru generarea de HTML, dar și pentru orice format text structurat (YAML, JSON, configurații).

În InfraMind, Jinja2 este utilizat pentru generarea fișierelor de configurare K3s (\texttt{config.yaml}) care sunt încărcate pe fiecare nod înainte de lansarea instalatorului. Există două template-uri:

\begin{itemize}
  \item \textbf{\texttt{master.yaml.j2}} — Template pentru nodurile master; conține logică condițională pentru nodul primordial (\texttt{cluster-init: true}) față de masterii secundari care se alătură clusterului existent (\texttt{server: "https://\{\{ primordial\_ip \}\}:6443"}). Include și configurații pentru TLS SAN, token, permisiuni kubeconfig și dezactivarea componentelor opționale (servicelb, traefik).
  \item \textbf{\texttt{worker.yaml.j2}} — Template pentru nodurile worker; conține adresa masterului primordial și token-ul de autentificare.
\end{itemize}

Prin separarea configurației într-un template Jinja2, InfraMind poate genera fișiere de configurare corecte și consistente pentru orice topologie de cluster, indiferent de numărul de noduri sau adresele IP ale acestora.

\section{Flask}

Flask este un micro-framework web scris în Python, conceput pentru a fi minimal și ușor extensibil [3]. Spre deosebire de framework-uri mai grele precum Django, Flask nu impune o structură rigidă, oferind flexibilitate maximă în organizarea aplicației.

În InfraMind, Flask îndeplinește simultan funcțiile de backend și de server pentru interfața web:

\begin{itemize}
  \item \textbf{Servirea interfeței utilizatorului} — Flask servește fișierul \texttt{index.html} și resursele statice (CSS, JavaScript), funcționând ca un server web simplu pentru aplicația single-page.
  \item \textbf{Endpoint-uri REST} — Backend-ul expune un set de rute HTTP pentru generarea inventarului, testarea conexiunilor SSH, lansarea și monitorizarea instalării, precum și interogarea clusterului prin \texttt{kubectl}.
  \item \textbf{Server-Sent Events (SSE)} — Pentru streaming-ul în timp real al progresului de instalare, Flask utilizează SSE (\texttt{text/event-stream}), permițând frontend-ului să primească actualizări fără polling activ.
  \item \textbf{Arhitectură Blueprint} — Aplicația este structurată în Blueprint-uri independente (\texttt{inventory\_bp}, \texttt{installer\_bp}, \texttt{uninstaller\_bp}, \texttt{kubectl\_bp}), fiecare responsabil de un domeniu funcțional separat. Fișierul \texttt{main.py} rămâne un app factory de 22 de linii.
\end{itemize}

\section{Python și librării asociate}

Întregul backend al InfraMind este implementat în Python 3 [14], beneficiind de ecosistemul bogat de librării disponibile:

\begin{itemize}
  \item \textbf{Paramiko} [11] — Librărie Python pentru comunicare SSH, utilizată ca motor principal de orchestrare. Permite executarea de comenzi la distanță, transferul de fișiere și testarea conectivității pe nodurile infrastructurii.
  \item \textbf{PyYAML} [13] — Utilizat pentru generarea și parsarea fișierelor de inventar în format YAML. Fiecare nod al clusterului este stocat ca un fișier YAML independent în directorul \texttt{inventory/}.
  \item \textbf{Jinja2} [12] — Utilizat pentru renderizarea template-urilor de configurare K3s (\texttt{master.yaml.j2}, \texttt{worker.yaml.j2}) cu parametrii specifici fiecărui nod.
  \item \textbf{threading / os / signal} — Module standard Python folosite pentru gestionarea stării proceselor de instalare, a flag-ului de anulare (\texttt{abort\_flag}) și a căilor de fișiere.
\end{itemize}

\section{JavaScript, HTML și CSS}

Interfața grafică a InfraMind este o aplicație web single-page (SPA) construită cu tehnologii web standard, fără a folosi un framework front-end complex:

\begin{itemize}
  \item \textbf{JavaScript vanilla (ES6+)} — Codul front-end este organizat în 10 module independente (\texttt{state.js}, \texttt{utils.js}, \texttt{navigation.js}, \texttt{inventory.js}, \texttt{topology.js}, \texttt{connections.js}, \texttt{deploy.js}, \texttt{cluster.js}, \texttt{handlers.js}, \texttt{main.js}), fiecare responsabil de o secțiune specifică a aplicației. Comunicarea cu backend-ul se face prin \texttt{fetch()} și \texttt{EventSource}.
  \item \textbf{HTML5} — Structura paginii este definită într-un singur fișier \texttt{index.html}, cu secțiuni ascunse/afișate dinamic prin JavaScript.
  \item \textbf{CSS3} — Stilizarea este împărțită în 9 fișiere CSS modulare (\texttt{base.css}, \texttt{buttons.css}, \texttt{navigation.css}, \texttt{inventory.css}, \texttt{topology.css}, \texttt{connections.css}, \texttt{deploy.css}, \texttt{welcome.css}, \texttt{cluster.css}), folosind variabile CSS custom (\texttt{:root vars}), animații (\texttt{@keyframes}) și layout flexibil.
  \item \textbf{Three.js / WebGL} [17] — Efectul vizual de fundal (\textit{light pillar}) este implementat cu Three.js, adăugând un element estetic modern interfeței.
\end{itemize}

\section{Bash}

Bash (\textit{Bourne Again Shell}) este shell-ul implicit al majorității distribuțiilor Linux și reprezintă limbajul de scripting utilizat de scripturile oficiale de instalare și dezinstalare ale K3s [18]. În contextul InfraMind, Bash intervine în două moduri principale:

\begin{itemize}
  \item \textbf{Scripturi de instalare K3s} — Fiecare nod master și worker rulează scriptul oficial \texttt{get.k3s.io} prin intermediul comenzii \texttt{curl -sfL https://get.k3s.io | sudo sh -s - server|agent}, lansată via SSH din Python. Scriptul descarcă și configurează K3s cu parametrii predefiniți în fișierul \texttt{/etc/rancher/k3s/config.yaml} generat anterior.
  \item \textbf{Scripturi de dezinstalare} — K3s furnizează scripturile \texttt{k3s-uninstall.sh} și \texttt{k3s-agent-uninstall.sh} care curăță complet instalarea de pe nodurile master, respectiv worker. Aceste scripturi sunt invocate din Python prin comenzi SSH executate cu Paramiko.
\end{itemize}

Deși nu este un limbaj de programare cu scop general, Bash este indispensabil în ecosistemul Linux pentru automatizarea sarcinilor de nivel sistem — un rol pe care îl îndeplinește eficient în tandem cu orchestrarea Python/SSH în cadrul InfraMind.

\section{Docker și Docker Compose}

Docker este platforma standard de containerizare care permite ambalarea unei aplicații împreună cu toate dependențele sale într-un container portabil [1]. InfraMind include un \texttt{Dockerfile} și un fișier \texttt{docker-compose.yml} pentru a simplifica rularea aplicației în orice mediu fără configurare manuală a dependențelor Python sau a serviciilor adiacente.

Utilizarea Docker în contextul InfraMind oferă:

\begin{itemize}
  \item \textbf{Portabilitate} — Aplicația rulează identic indiferent de sistemul de operare al gazdei.
  \item \textbf{Izolare} — Dependențele Python (Flask, Paramiko, PyYAML, Jinja2) sunt izolate în container, evitând conflictele cu pachetele sistemului.
  \item \textbf{Deployment simplificat} — O singură comandă (\texttt{docker compose up}) pornește întreaga aplicație.
\end{itemize}

\chapter{Cerințele aplicației}
\label{cap:cerinte}
\section{Scop}

Scopul principal al InfraMind este de a elimina complexitatea asociată cu instalarea și administrarea Kubernetes, oferind o experiență fluidă și accesibilă chiar și pentru cei care nu sunt experți în orchestrarea containerelor. Aplicația se adresează atât administratorilor de infrastructură, cât și dezvoltatorilor care doresc să implementeze rapid un cluster Kubernetes fără a pierde timp cu configurări manuale și procese complicate.

Prin automatizarea implementării clusterelor K3s și integrarea cu diverse tipuri de infrastructură, InfraMind contribuie la eficientizarea operațiunilor IT, reducerea timpului necesar pentru configurare și îmbunătățirea scalabilității resurselor. În plus, interfața web ușor de utilizat permite utilizatorilor să monitorizeze și să gestioneze clusterele într-un mod centralizat, oferind un control sporit asupra întregului ecosistem Kubernetes.

\section{User experience}
Experiența utilizatorului este esențială pentru relevanța soluției noastre, ea trebuie sa fie una cat mai plăcută, fluidă și simplă, astfel încât utilizatorii să poată instala și gestiona clustere K3s fără a fi necesare cunoștințe avansate despre Kubernetes. Interfața web trebuie să fie intuitivă, cu un design clar și elemente vizuale care să ghideze utilizatorul pas cu pas prin procesul de instalare.

Navigarea trebuie să fie logică și bine structurată, evitând pași inutili sau formulare complicate. De asemenea, trebuie să oferim feedback vizual clar pentru fiecare acțiune efectuată, cum ar fi bare de progres, notificări și mesaje de confirmare.

Pe lângă simplitate, aplicația trebuie să asigure o experiență fluidă, minimizând timpii de așteptare și eliminând orice blocaj care ar putea încetini procesul. Automatismele și presetările inteligente vor ajuta utilizatorii să finalizeze instalarea rapid, reducând necesitatea unor configurări manuale complexe.

În plus, integrarea cu infrastructuri diverse trebuie să fie transparentă pentru utilizator, iar erorile sau problemele întâlnite să fie explicate clar, oferind soluții sau pași de remediere direct în interfață. Astfel, InfraMind nu doar că facilitează implementarea clusterelor K3s, dar creează și o experiență plăcută, eliminând barierele tehnice pentru utilizatori.

Pentru a menține această experiență cât mai accesibilă, InfraMind trebuie să includă un proces de onboarding bine structurat, astfel încât utilizatorii noi să poată începe rapid fără a fi copleșiți de setări complexe. Acest lucru poate fi realizat prin tutoriale interactive, sugestii contextuale și setări implicite optimizate pentru cele mai comune scenarii de utilizare.

De asemenea, personalizarea interfeței joacă un rol important în menținerea unui flux de lucru eficient. Utilizatorii ar trebui să aibă posibilitatea de a ajusta configurațiile în funcție de nevoile lor, fără a fi forțați să navigheze prin meniuri complexe. Opțiuni precum presetări recomandate, profiluri de configurare și automatizări pentru instalarea clusterelor ar contribui semnificativ la îmbunătățirea experienței.

Un alt aspect esențial este responsivitatea interfeței, asigurându-ne că aplicația funcționează optim pe orice dispozitiv, fie că este vorba de un laptop sau desktop. De asemenea, performanța trebuie optimizată astfel încât acțiunile utilizatorilor să fie procesate rapid, fără întârzieri sau blocaje.

În plus, suportul și documentația trebuie să fie integrate direct în aplicație, astfel încât utilizatorii să aibă acces rapid la răspunsuri și soluții în cazul în care întâmpină dificultăți. Un sistem de notificări și alerte în timp real poate ajuta utilizatorii să rămână informați despre statusul clusterelor lor și eventualele probleme.

Prin implementarea acestor principii, InfraMind va reuși să ofere o experiență de utilizare fluidă, accesibilă și eficientă, permițând oricui să instaleze și să administreze clustere K3s fără efort inutil.

\section{Eficiență}

Eficiența este un pilon central al InfraMind, atât în ceea ce privește performanța aplicației, cât și procesul de instalare și gestionare a clusterelor K3s. Scopul este reducerea timpului necesar pentru implementare, eliminând pașii inutili și optimizând fiecare etapă astfel încât utilizatorii să poată avea un cluster funcțional în cel mai scurt timp posibil.

\subsection{Automatizarea procesului de instalare}

Unul dintre cele mai mari obstacole în implementarea unui cluster Kubernetes tradițional este necesitatea unei configurații manuale extinse. InfraMind elimină această barieră printr-un proces de instalare automatizat, care:

\begin{itemize}
  \item Detectează și validează infrastructura disponibilă (fizică sau virtuală).
  \item Configurează nodurile, rolurile și resursele clusterului fără intervenția utilizatorului.
  \item Aplică setări optimizate pentru performanță, fără a necesita ajustări manuale complexe.
\end{itemize}

Astfel, utilizatorii pot avea un cluster K3s operațional în câteva minute, comparativ cu ore sau chiar zile în cazul unei configurări manuale.

\subsection{Optimizarea resurselor}

Un alt aspect cheie al eficienței este utilizarea optimă a resurselor hardware și software. InfraMind asigură acest lucru prin:

\begin{itemize}
  \item Alocarea dinamică a resurselor, evitând supra-provizionarea sau sub-provizionarea clusterului.
  \item Detectarea și corectarea automată a erorilor comune, reducând timpul petrecut pe depanare.
  \item Integrarea cu infrastructuri diverse, permițând rularea eficientă pe hardware low-power sau în medii cloud cost-eficiente.
\end{itemize}

Prin aceste optimizări, aplicația garantează un consum minim de resurse, maximizând performanța clusterului fără risipă de procesare, memorie sau stocare.

\subsection{Reducerea complexității operaționale}

Gestionarea unui cluster Kubernetes poate fi un proces complex și predispus la erori. InfraMind rezolvă această problemă printr-un set de funcționalități care simplifică operațiunile de zi cu zi:

\begin{itemize}
  \item Feedback vizual și notificări în timp real, oferind utilizatorilor informații clare despre starea clusterului.
  \item Mentenanță simplificată prin automatizarea pașilor repetitivi de configurare.
  \item Suport pentru scalare, adaptând configurația clusterului în funcție de necesități.
\end{itemize}

Aceste caracteristici permit utilizatorilor să administreze clusterele cu un efort minim, reducând intervențiile manuale și optimizând fluxurile de lucru.

\subsection{Performanță optimizată a aplicației}

Pentru a asigura o experiență fluidă și rapidă, InfraMind este construit pe principii moderne de dezvoltare software:

\begin{itemize}
  \item Arhitectură scalabilă, capabilă să gestioneze multiple instalări simultan fără degradarea performanței.
  \item Interfață web rapidă, optimizată pentru timpi de răspuns minimali și interacțiuni eficiente.
  \item Procesare asincronă prin Server-Sent Events, astfel încât utilizatorii să primească feedback în timp real fără să blocheze interacțiunea cu aplicația.
\end{itemize}

Toate aceste aspecte contribuie la un proces de instalare și administrare care este nu doar rapid, ci și stabil și fiabil, permițând utilizatorilor să se concentreze pe rularea aplicațiilor lor în Kubernetes, fără griji legate de infrastructură.

\subsection{Eficiență pe termen lung}

Pe lângă instalare, InfraMind se axează și pe menținerea eficienței în timp, prin:

\begin{itemize}
  \item Compatibilitate extinsă, suportând diverse tipuri de infrastructură și configurări personalizate.
  \item Integrare cu Docker și cu uneltele standard ale ecosistemului Python, permițând extinderea și personalizarea fluxurilor de lucru.
  \item Posibilitatea de a refolosi inventarele și configurațiile generate pentru noi instalări.
\end{itemize}

Astfel, aplicația nu doar că optimizează implementarea inițială, ci asigură și un management eficient al clusterelor pe termen lung, reducând costurile și efortul operațional al utilizatorilor.

\chapter{Structura aplicației}
\label{cap:structura}
\section{Overview}

InfraMind este o aplicație web Python/Flask care reunește în același proces atât frontend-ul (interfața utilizatorului) cât și backend-ul (orchestrarea instalării). Spre deosebire de arhitecturile tradiționale care delegă instalarea unui tool extern de tip configuration management, InfraMind implementează propriul motor de orchestrare SSH direct în Python, fără dependențe la runtime în afara celor din \texttt{requirements.txt}.

Arhitectura aplicației urmărește principiul separării responsabilităților: fiecare modul Python gestionează un domeniu bine definit, iar comunicarea între frontend și backend se realizează prin apeluri HTTP REST și prin fluxuri SSE pentru actualizări în timp real. Această relație se poate observa în \hyperref[fig01.png]{Figura~\ref{fig01.png}}.

\begin{figure}[H]
    \centering
    \includegraphics[width=1\textwidth]{fig01.png}
    \caption{Overview general al aplicației InfraMind}
    \label{fig01.png}
\end{figure}

\section{Fluxul de lucru al aplicației}

Utilizatorul parcurge un proces ghidat în patru pași prin intermediul interfeței web:

\begin{enumerate}
  \item \textbf{Definirea inventarului} — Utilizatorul adaugă nodurile clusterului (adresă IP, rol master/worker) și desemnează nodul master primordial. Backend-ul scrie câte un fișier YAML per nod în directorul \texttt{inventory/}.
  \item \textbf{Validarea topologiei} — Tab-ul Topology afișează vizual structura clusterului, permițând utilizatorului să verifice configurația înainte de continuare.
  \item \textbf{Testarea conexiunilor SSH} — Tab-ul Connections permite testarea accesibilității SSH pentru fiecare nod individual, utilizând cheia privată SSH furnizată de utilizator.
  \item \textbf{Instalarea clusterului} — Tab-ul Deploy primește acreditările SSH (username și cheie privată), lansează procesul de instalare și afișează progresul în timp real prin Server-Sent Events.
\end{enumerate}

Pe lângă fluxul de instalare, aplicația oferă și o secțiune separată pentru gestionarea unui cluster existent, accesibilă prin încărcarea unui fișier kubeconfig.

\section{Structuri interne}

\subsection{Frontend — Interfața utilizatorului}

Frontend-ul este o aplicație Single-Page construită cu JavaScript vanilla (ES6+), servită direct de Flask din directorul \texttt{frontend/src/templates/}. Interfața este organizată în două fluxuri principale: fluxul de instalare cluster nou și panoul de gestionare cluster existent.

Fluxul de instalare este structurat în patru tab-uri cu navigare pas cu pas:

\begin{itemize}
  \item \textbf{Inventory} — formular pentru adăugarea și editarea nodurilor; include un selector de rol (master/worker), câmp pentru adresa IP și selector de nod primordial.
  \item \textbf{Topology} — vizualizare SVG a topologiei clusterului cu editare inline; nodurile sunt reprezentate grafic împreună cu conexiunile dintre ele.
  \item \textbf{Connections} — panou de testare SSH cu indicatori vizuali per nod (verde/roșu/în progres) și câmp pentru introducerea cheii SSH private.
  \item \textbf{Deploy} — panoul de instalare cu carduri de progres animate, buton de anulare și panoul de afișare a kubeconfig-ului după instalare reușită.
\end{itemize}

Codul JavaScript este organizat în 10 module independente, fiecare responsabil de o zonă funcțională, conform tabelului următor:

\begin{center}
\begin{tabular}{|l|l|}
\hline
\textbf{Modul} & \textbf{Responsabilitate} \\
\hline
\texttt{state.js}      & Toate variabilele globale partajate \\
\texttt{utils.js}      & \texttt{escapeHtml}, \texttt{showToast}, \texttt{showConfirmToast} \\
\texttt{navigation.js} & Comutarea tab-urilor, animație bubble indicator \\
\texttt{inventory.js}  & Gestionarea listei de noduri, \texttt{renderVMList}, \texttt{detectInventory} \\
\texttt{topology.js}   & Vizualizarea topologiei, editorul SVG inline \\
\texttt{connections.js}& Testarea conexiunilor SSH per nod \\
\texttt{deploy.js}     & Step cards animate, \texttt{startDeploy}, \texttt{abortDeploy}, \texttt{startUninstall} \\
\texttt{cluster.js}    & Welcome screen, dashboard cluster, carduri de resurse per nod \\
\texttt{handlers.js}   & \texttt{setupHandlers()} — toate legăturile de eveniment \\
\texttt{main.js}       & Blocul \texttt{DOMContentLoaded} de inițializare \\
\hline
\end{tabular}
\end{center}

Stilizarea este împărțită în 9 fișiere CSS modulare: \texttt{base.css}, \texttt{buttons.css}, \texttt{navigation.css}, \texttt{inventory.css}, \texttt{topology.css}, \texttt{connections.css}, \texttt{deploy.css}, \texttt{welcome.css}, \texttt{cluster.css}. Fiecare fișier gestionează exclusiv stilurile componentei sale, folosind variabile CSS custom definite în \texttt{:root} din \texttt{base.css}.

Comunicarea cu backend-ul se realizează prin două mecanisme:
\begin{itemize}
  \item \texttt{fetch()} — pentru toate apelurile REST sincrone (generate inventar, testare SSH, kubectl).
  \item \texttt{EventSource} — pentru fluxurile SSE de progres în timp real (instalare, dezinstalare).
\end{itemize}

Aspectul general al frontend-ului este ilustrat în \hyperref[frontend.png]{Figura~\ref{frontend.png}}.

\begin{figure}[H]
    \centering
    \includegraphics[width=1\textwidth]{frontend.png}
    \caption{Interfața web a aplicației InfraMind}
    \label{frontend.png}
\end{figure}

\subsection{Backend — Arhitectura modulară}

Backend-ul Flask este implementat ca un app factory thin (\texttt{main.py}, 22 de linii) care importă și înregistrează patru Blueprint-uri. Structura modulară a fost aleasă pentru a menține o separare clară a responsabilităților și pentru a facilita extinderea ulterioară a aplicației.

\begin{center}
\begin{tabular}{|l|l|l|}
\hline
\textbf{Fișier} & \textbf{Blueprint} & \textbf{Responsabilitate} \\
\hline
\texttt{main.py}        & —                  & App factory; înregistrare Blueprint-uri \\
\texttt{config.py}      & —                  & Constante de cale, stare partajată, \texttt{abort\_flag} \\
\texttt{ssh.py}         & —                  & Funcții auxiliare SSH (Paramiko) \\
\texttt{inventory.py}   & \texttt{inventory\_bp}  & CRUD inventar, testare SSH \\
\texttt{installer.py}   & \texttt{installer\_bp}  & Instalare K3s via SSH \\
\texttt{uninstaller.py} & \texttt{uninstaller\_bp}& Dezinstalare K3s via SSH \\
\texttt{kubectl.py}     & \texttt{kubectl\_bp}    & Operații kubectl pe cluster existent \\
\hline
\end{tabular}
\end{center}

Modulul \texttt{config.py} definește starea partajată a aplicației: un obiect \texttt{\_DeployState} cu câmpul \texttt{status} (valorile posibile: \texttt{idle}, \texttt{running}, \texttt{success}, \texttt{failed}, \texttt{aborted}), un \texttt{threading.Lock} (\texttt{proc\_lock}) pentru accesul exclusiv la procesele de instalare/dezinstalare și un \texttt{threading.Event} (\texttt{abort\_flag}) pentru semnalarea anulării.

Lista completă a endpoint-urilor HTTP expuse de aplicație este prezentată în tabelul următor:

\begin{center}
\begin{tabular}{|l|l|l|}
\hline
\textbf{Metodă} & \textbf{Rută} & \textbf{Scop} \\
\hline
GET  & \texttt{/}                      & Servire \texttt{index.html} \\
POST & \texttt{/generate}              & Scriere fișiere YAML inventar \\
GET  & \texttt{/detect-inventory}      & Scanare director \texttt{inventory/} \\
POST & \texttt{/delete-host}           & Ștergere YAML nod \\
POST & \texttt{/test-ssh}              & Testare conectivitate SSH per nod \\
GET  & \texttt{/deploy}                & SSE stream — instalare K3s \\
POST & \texttt{/deploy-abort}          & Setare \texttt{abort\_flag} \\
GET  & \texttt{/deploy-status}         & Stare curentă (idle/running/…) \\
GET  & \texttt{/uninstall}             & SSE stream — dezinstalare K3s \\
POST & \texttt{/kubectl-nodes}         & \texttt{kubectl get nodes -o wide} \\
POST & \texttt{/kubectl-pods}          & \texttt{kubectl get pods -A -o wide} \\
POST & \texttt{/kubectl-services}      & \texttt{kubectl get services -A} \\
POST & \texttt{/kubectl-node-resources}& Metrici CPU/memorie per nod \\
\hline
\end{tabular}
\end{center}

\subsection{Motorul de orchestrare SSH}

Modulul \texttt{ssh.py} constituie fundația motorului de orchestrare și expune trei funcții utilizate de \texttt{installer.py} și \texttt{uninstaller.py}:

\begin{itemize}
  \item \textbf{\texttt{\_write\_temp\_key(ssh\_key\_text)}} — Scrie cheia SSH privată furnizată de utilizator într-un fișier temporar cu permisiuni stricte (\texttt{0600}) și returnează calea acestuia. Fișierul este șters automat după finalizarea operației.
  \item \textbf{\texttt{\_open\_ssh\_client(ip, username, key\_path)}} — Creează și conectează un client SSH Paramiko. Încearcă succesiv formatele de cheie Ed25519, RSA, ECDSA și DSS, asigurând compatibilitatea cu orice infrastructură existentă.
  \item \textbf{\texttt{\_ssh\_run\_live(client, cmd)}} — Generator non-blocant care yield-uiește \texttt{(linie, None)} pentru fiecare linie de output a comenzii și \texttt{(None, exit\_code)} la finalizare. Canalul SSH operează fără blocare, iar \texttt{abort\_flag} este verificat la fiecare iterație pentru a permite anularea imediată a instalării.
\end{itemize}

\subsection{Gestionarea inventarului}

Fiecare nod al clusterului este reprezentat printr-un fișier YAML independent stocat în directorul \texttt{inventory/} din rădăcina workspace-ului. Schema unui fișier de inventar este minimală:

\begin{verbatim}
name: master1
ip: 192.168.1.10
role: master
primordial: true   # prezent doar la nodul master primordial
\end{verbatim}

Separarea fiecărui nod într-un fișier distinct permite operații granulare (adăugare/ștergere individuală) și evită conflictele de concurență la scriere. La detectarea inventarului existent, backend-ul scanează directorul și reconstituie lista de noduri în ordinea alfabetică a numelor de fișiere, identificând automat nodul primordial.

\subsection{Instalarea clusterului K3s prin SSH}

Procesul de instalare este implementat în \texttt{installer.py} ca un generator Python (\texttt{\_stream\_k3s\_install}) care yield-uiește evenimente SSE pe măsură ce fiecare pas avansează. Ordinea de execuție este strictă și motivată de dependențele tehnice ale K3s:

\begin{enumerate}
  \item \textbf{Instalarea Docker (opțional)} — Pe fiecare nod se verifică prezența Docker (\texttt{docker --version}). Dacă Docker nu este instalat și utilizatorul a bifat opțiunea corespunzătoare, se rulează scriptul oficial \texttt{https://get.docker.com} prin \texttt{curl | sudo sh}.
  \item \textbf{Inițializarea masterului primordial} — Se renderizează template-ul Jinja2 \texttt{master.yaml.j2} cu parametrii \texttt{cluster\_init: true}, token-ul K3s, adresa IP și flag-ul Docker. Fișierul rezultat este încărcat pe nod la \texttt{/etc/rancher/k3s/config.yaml} prin \texttt{sudo tee}, după care se rulează \texttt{curl -sfL https://get.k3s.io | sudo sh -s - server}.
  \item \textbf{Extragerea kubeconfig} — Backend-ul descarcă \texttt{/etc/rancher/k3s/k3s.yaml} de pe nodul primordial, înlocuiește adresa \texttt{127.0.0.1} cu IP-ul extern al masterului și salvează fișierul local la \texttt{\textasciitilde/.kube/k3s.yaml}.
  \item \textbf{Alăturarea masterilor adicionali} — Pentru fiecare master secundar, se renderizează \texttt{master.yaml.j2} cu \texttt{cluster\_init: false} și adresa primordialului, după care se execută același script de instalare K3s.
  \item \textbf{Înrolarea nodurilor worker} — Se renderizează \texttt{worker.yaml.j2} cu adresa masterului și token-ul, iar scriptul de instalare K3s este rulat cu argumentul \texttt{agent}.
\end{enumerate}

La finalizarea cu succes a tuturor pașilor, evenimentul SSE \texttt{finished} include și conținutul complet al kubeconfig-ului, permițând frontend-ului să îl afișeze imediat utilizatorului.

\subsection{Dezinstalarea clusterului}

Dezinstalarea este implementată simetric față de instalare, în ordine inversă: mai întâi workerii, apoi masterii adicionali, în final masterul primordial. Pe fiecare nod se execută secvențial:

\begin{itemize}
  \item \texttt{sudo /usr/local/bin/k3s-agent-uninstall.sh} (pe workeri) sau \texttt{sudo /usr/local/bin/k3s-uninstall.sh} (pe masteri).
  \item \texttt{docker kill \$(docker ps -q) || true} și \texttt{docker rm -f \$(docker ps -aq) || true} pentru curățarea tuturor containerelor Docker rulate de K3s. Comenzile sunt sufixate cu \texttt{|| true} pentru a nu eșua în absența Docker sau a containerelor.
\end{itemize}

La final, fișierul local de kubeconfig \texttt{\textasciitilde/.kube/k3s.yaml} este șters. Întregul proces este transmis frontend-ului în timp real prin același mecanism SSE folosit la instalare.

\subsection{Mecanismul de anulare (abort)}

Utilizatorul poate anula o instalare sau dezinstalare aflată în desfășurare prin butonul Abort din interfață. La apăsare, frontend-ul trimite un \texttt{POST /deploy-abort}, care setează \texttt{abort\_flag} (un \texttt{threading.Event} Python). La iterația următoare a generatorului SSH (\texttt{\_ssh\_run\_live}), flag-ul este detectat, canalul SSH curent este închis, iar generatorul returnează fără a mai executa pașii următori. Starea se actualizează la \texttt{aborted} și frontend-ul afișează butonul de reluare.

\subsection{Dashboard-ul de cluster și monitorizarea resurselor}

Secțiunea \textit{Existing Cluster} permite utilizatorilor să monitorizeze un cluster K3s existent prin încărcarea unui kubeconfig. Frontend-ul trimite kubeconfig-ul ca parte din payload-ul fiecărui apel kubectl. Backend-ul (\texttt{kubectl.py}) execută comanda kubectl corespunzătoare printr-un subprocess Python și returnează datele parsate.

Sunt disponibile patru vizualizări:

\begin{itemize}
  \item \textbf{Nodes} — tabelul nodurilor cu roluri, versiunea K3s și statusul fiecăruia.
  \item \textbf{Pods} — tabelul pod-urilor din toate namespace-urile, cu starea și nodul pe care rulează.
  \item \textbf{Services} — tabelul serviciilor cu tipul, porturile expuse și clusterIP.
  \item \textbf{Resources} — carduri per nod cu bare de progres pentru CPU, memorie și număr de pod-uri. Datele de bază provin din \texttt{kubectl get nodes -o json} (capacitate alocabilă), iar datele de utilizare în timp real provin din \texttt{kubectl top nodes} (disponibile dacă metrics-server este instalat). Culoarea barelor indică gradul de utilizare: verde sub 70\%, galben sub 90\%, roșu la 90\% și peste.
\end{itemize}

\subsection{Afișarea kubeconfig după instalare}

La finalizarea cu succes a unui deployment, evenimentul SSE \texttt{finished} include câmpul \texttt{kubeconfig} cu conținutul complet al fișierului \texttt{k3s.yaml}, cu adresa IP a masterului primordial deja introdusă. Frontend-ul afișează automat un panou colaibil cu kubeconfig-ul, formatat monospaciat pe fundal întunecat. Un buton de copiere în clipboard permite utilizatorului să preia imediat kubeconfig-ul pentru utilizare cu \texttt{kubectl} local.

Interacțiunea generală dintre componentele aplicației este ilustrată în \hyperref[overalluml.png]{Figura~\ref{overalluml.png}}.

\begin{figure}[H]
    \centering
    \includegraphics[height=0.5\textwidth]{overalluml.png}
    \caption{Diagrama de comunicare dintre componente}
    \label{overalluml.png}
\end{figure}

\chapter*{Concluzii}
\addcontentsline{toc}{chapter}{Concluzii}

InfraMind a fost conceput ca răspuns la o problemă concretă: deși Kubernetes a devenit standardul de facto pentru orchestrarea aplicațiilor containerizate, procesul de instalare și configurare a unui cluster rămâne complex și necesită cunoștințe tehnice avansate din domenii multiple. Lucrarea de față a documentat proiectarea și implementarea unei soluții care adresează această problemă printr-o interfață web intuitivă și un motor de orchestrare SSH implementat integral în Python.

\section*{Ce s-a realizat}

InfraMind oferă un flux complet de la zero la cluster funcțional, fără nicio dependență externă de tip configuration management pe mașina de control. Utilizatorul definește inventarul nodurilor, testează conexiunile SSH, lansează instalarea și primește kubeconfig-ul gata de utilizat — totul din browser, fără a scrie niciun fișier de configurare manual.

Principalele realizări tehnice ale proiectului sunt:

\begin{itemize}
  \item \textbf{Motor de orchestrare SSH pur Python} — Folosind Paramiko și generatoare Python, aplicația execută secvențe complexe de comenzi SSH pe mai multe noduri simultan, cu streaming live al output-ului și suport pentru anulare imediată.
  \item \textbf{Arhitectură modulară Flask cu Blueprint-uri} — Backend-ul este structurat în șapte module cu responsabilități clare, fără dependențe circulare, facilitând mentenanța și extinderea.
  \item \textbf{Frontend SPA modular} — Zece module JavaScript și nouă fișiere CSS independente formează o interfață single-page fluidă, cu animații, feedback vizual în timp real și navigare pas cu pas.
  \item \textbf{Server-Sent Events pentru progres în timp real} — Utilizatorul vede fiecare pas al instalării pe măsură ce se execută, cu stări vizuale clare (în așteptare, activ, finalizat, eșuat, anulat).
  \item \textbf{Dashboard de monitorizare cluster} — Secțiunea de cluster existent permite vizualizarea nodurilor, pod-urilor, serviciilor și metricilor de resurse (CPU, memorie, pod-uri) prin \texttt{kubectl}.
  \item \textbf{Suport pentru clustere HA} — Topologia cu mai mulți masteri (embedded etcd) este suportată nativ, cu inițializare corectă a masterului primordial și alăturare automată a masterilor secundari.
  \item \textbf{Containerizare Docker} — Aplicația este complet containerizată, putând fi lansată cu \texttt{docker compose up} în orice mediu fără configurare suplimentară.
\end{itemize}

\section*{Provocări întâlnite}

Pe parcursul implementării au apărut mai multe provocări tehnice interesante. Gestionarea output-ului SSH în timp real pe canale non-blocante a necesitat o abordare atentă pentru a evita blocarea thread-ului Flask și pentru a menține compatibilitatea cu mecanismul SSE. Mecanismul de abort a trebuit să fie fiabil la nivel de thread, folosind un \texttt{threading.Event} verificat la fiecare iterație a generatorului SSH. Gestionarea erorilor parțiale de instalare (un nod eșuează, ceilalți continuă) a necesitat o logică de fallthrough atentă pentru a asigura că utilizatorul primește feedback complet indiferent de scenariul de eroare.

O altă provocare a reprezentat-o tranziția arhitecturală de la o abordare bazată pe Ansible la orchestrare pură Python/SSH: toate abstractizările oferite de Ansible (inventar, roluri, handlere, idempotență) au trebuit reimplementate explicit în cod Python, ceea ce a condus la o arhitectură mai transparentă și mai ușor de înțeles, dar și mai verbosă.

\section*{Limitări}

Ca orice soluție orientată spre simplitate, InfraMind are și limitări față de soluțiile enterprise:

\begin{itemize}
  \item \textbf{O singură instalare activă simultan} — Mecanismul de stare partajat permite un singur proces de instalare/dezinstalare la un moment dat; suportul pentru instalări paralele ar necesita o arhitectură cu cozi de task-uri.
  \item \textbf{Fără autentificare} — Interfața web nu implementează autentificare; se presupune că aplicația rulează într-un mediu privat (rețea locală sau container izolat).
  \item \textbf{Fără persistență între sesiuni} — Starea deployment-ului (dacă clusterul este instalat sau nu) nu este persistată; la repornirea aplicației, utilizatorul trebuie să re-detecteze inventarul.
  \item \textbf{Volum de lucru mic spre mediu} — InfraMind este optimizat pentru clustere de dimensiune redusă spre medie (tipic 1--9 noduri), similar K3s față de Kubernetes standard.
\end{itemize}

\section*{Direcții viitoare}

Proiectul InfraMind poate fi extins în mai multe direcții naturale:

\begin{itemize}
  \item \textbf{Autentificare și autorizare} — Adăugarea unui sistem de login simplu (OAuth2 sau JWT) ar face aplicația sigur de expus pe internet.
  \item \textbf{Gestionarea mai multor clustere} — Suportul pentru profiluri de cluster (multiple inventare salvate și comutabile) ar extinde utilitatea aplicației.
  \item \textbf{Integrare Helm} — Un panou de deploy de chart-uri Helm direct din interfață ar completa ciclul de viață al clusterului.
  \item \textbf{Monitorizare integrată} — Instalarea automată a unui stack de monitorizare (Prometheus + Grafana) ca pas opțional post-instalare ar oferi vizibilitate completă.
  \item \textbf{Notificări} — Integrarea cu sisteme de notificare (email, Slack, webhook) la finalizarea/eșuarea instalărilor ar îmbunătăți experiența în medii de echipă.
\end{itemize}

În ansamblu, InfraMind demonstrează că automatizarea instalării Kubernetes nu necesită neapărat tooling enterprise complex — o aplicație Python bine structurată, cu orchestrare SSH directă și o interfață web modernă, poate oferi o experiență de utilizare superioară pentru segmentul de utilizatori vizat: administratori de infrastructură și dezvoltatori care doresc să implementeze rapid clustere K3s pe orice tip de hardware.

\chapter*{Rezumat}
\addcontentsline{toc}{chapter}{Rezumat}

\section*{Rezumat în limba română}

InfraMind este o aplicație web open-source pentru automatizarea instalării și gestionării clusterelor Kubernetes K3s. Implementată în Python cu Flask, aplicația oferă o interfață web intuitivă prin care utilizatorii pot defini arhitectura clusterului (noduri master și worker), testa conexiunile SSH, lansa instalarea și monitoriza progresul în timp real.

Orchestrarea instalării se realizează direct prin SSH, folosind librăria Paramiko, fără dependențe externe de tip configuration management. Configurațiile K3s per nod sunt generate dinamic prin template-uri Jinja2, asigurând corectitudine și consistență indiferent de topologia aleasă. Progresul instalării este transmis frontend-ului în timp real prin Server-Sent Events, oferind utilizatorului feedback vizual detaliat la fiecare pas.

Aplicația include și un dashboard pentru gestionarea clusterelor existente, cu vizualizare de noduri, pod-uri, servicii și metrici de resurse (CPU, memorie) obținute prin kubectl. La finalizarea cu succes a instalării, kubeconfig-ul este afișat automat și poate fi copiat cu un singur click.

Aplicația este containerizată cu Docker și poate fi pornită cu o singură comandă. Codul sursă este organizat modular atât în backend (Blueprint-uri Flask) cât și în frontend (module JavaScript și CSS independente), facilitând mentenanța și extinderea.

\section*{Abstract}

InfraMind is an open-source web application for automating the installation and management of K3s Kubernetes clusters. Built in Python using the Flask framework, it provides an intuitive web interface for defining cluster architecture (master and worker nodes), testing SSH connectivity, launching installations, and monitoring progress in real time.

The installation orchestration runs entirely over SSH using the Paramiko library, with no external configuration management dependencies such as Ansible. Per-node K3s configurations are generated dynamically from Jinja2 templates, ensuring correctness and consistency regardless of the chosen topology. Installation progress is streamed to the frontend in real time via Server-Sent Events, providing the user with detailed visual feedback at each step.

The application also includes a cluster management dashboard for visualising nodes, pods, services, and resource metrics (CPU, memory) retrieved via kubectl. Upon successful deployment, the kubeconfig is automatically displayed and can be copied to the clipboard with a single click.

InfraMind is fully containerised with Docker and can be started with a single \texttt{docker compose up} command. The codebase is organised modularly in both backend (Flask Blueprints) and frontend (independent JavaScript and CSS modules), facilitating maintenance and future extension.



\chapter*{Bibliografie}
\addcontentsline{toc}{chapter}{Bibliografie}

\begin{enumerate}

\item Docker Inc., \textit{Docker Documentation}, disponibil la: \url{https://docs.docker.com/}, accesat 30 martie 2026.

\item Edge Delta, \textit{State of Kubernetes 2023 — Adoption Trends}, disponibil la: \url{https://edgedelta.com/company/blog/state-of-kubernetes-2023}, accesat 30 martie 2026.

\item Flask (Pallets Projects), \textit{Flask Documentation}, disponibil la: \url{https://flask.palletsprojects.com/en/stable/}, accesat 30 martie 2026.

\item Google Cloud, \textit{Google Kubernetes Engine Documentation}, disponibil la: \url{https://cloud.google.com/kubernetes-engine/docs}, accesat 30 martie 2026.

\item K3s (SUSE/Rancher Labs), \textit{K3s — Lightweight Kubernetes}, disponibil la: \url{https://k3s.io/}, accesat 30 martie 2026.

\item Kind (Kubernetes SIGs), \textit{Kind Documentation — Kubernetes IN Docker}, disponibil la: \url{https://kind.sigs.k8s.io/}, accesat 30 martie 2026.

\item Kubernetes, \textit{Kubeadm Documentation}, disponibil la: \url{https://kubernetes.io/docs/reference/setup-tools/kubeadm/}, accesat 30 martie 2026.

\item Kubernetes, \textit{Kubernetes Documentation}, disponibil la: \url{https://kubernetes.io/docs/}, accesat 30 martie 2026.

\item Microsoft, \textit{Azure Kubernetes Service Documentation}, disponibil la: \url{https://learn.microsoft.com/en-us/azure/aks/}, accesat 30 martie 2026.

\item Minikube (Kubernetes SIGs), \textit{Minikube Documentation}, disponibil la: \url{https://minikube.sigs.k8s.io/docs/}, accesat 30 martie 2026.

\item Paramiko, \textit{Paramiko Documentation — SSH2 for Python}, disponibil la: \url{https://www.paramiko.org/}, accesat 30 martie 2026.

\item Pallets Projects, \textit{Jinja2 Documentation — Template Engine for Python}, disponibil la: \url{https://jinja.palletsprojects.com/}, accesat 30 martie 2026.

\item PyYAML, \textit{PyYAML Documentation}, disponibil la: \url{https://pyyaml.org/wiki/PyYAMLDocumentation}, accesat 30 martie 2026.

\item Python Software Foundation, \textit{Python 3 Documentation}, disponibil la: \url{https://docs.python.org/3/}, accesat 30 martie 2026.

\item Rancher (SUSE), \textit{Rancher Documentation}, disponibil la: \url{https://www.rancher.com/docs/}, accesat 30 martie 2026.

\item Red Hat, \textit{OpenShift Documentation}, disponibil la: \url{https://docs.openshift.com/}, accesat 30 martie 2026.

\item Three.js, \textit{Three.js Documentation}, disponibil la: \url{https://threejs.org/docs/}, accesat 30 martie 2026.

\item GNU Project, \textit{Bash Reference Manual}, disponibil la: \url{https://www.gnu.org/software/bash/manual/}, accesat 30 martie 2026.

\end{enumerate}

\end{document}